\begin{python0}
from solutions import *; clear()
\end{python0}
\sectionthree{Type conversion}

Add the following to our Basic Example:
\begin{console}[commandchars=\~\@\$]
#include <iostream>

class P
{     
      ...
};

class C: public P
{
        ...
};

int main()
{   
    C c;
    ...
    P p;
    ...
    
    ~textbf@C c0(p); // This does not work!!!$
    ~textbf@P p0(c); // This works$

    return 0;
}
\end{console}

This tells you that if \verb!C! is a subclass of \verb!P!, then automatically, \verb!P! has a constructor of the form
\begin{console}
class P
{     
      ...
      P(const C &);
      ...
};
\end{console}

The default behavior of this constructor is to copy relevant values from
a \verb!C! object to the \verb!P! object that is being initialized.

This acts as a type conversion operator.

The idea is very similar to the assignment operator between objects of
\verb!C! and objects or \verb!P!. And the reason that
\begin{console}
C c0(p);
\end{console}

does NOT work is similar.

